\annex{обязательное}{Справка об исследовании патентной литературы}
\label{sec:appendix_a}

Наименование объекта поиска: Гибридный алгоритм сжатия изображений на основе вейвлет-преобразования и нейросетевого восстановления сверхразрешения (ESPCN, FSRCNN).
Что и за какой период просмотрено:

Наименование объекта поиска: Гибридный алгоритм сжатия изображений на основе вейвлет-преобразования и нейросетевого восстановления сверхразрешения (ESPCN, FSRCNN).
Что и за какой период просмотрено:
\begin{itemize}
    \item электронная база данных Национального центра интеллектуальной собственности --- \url{http://www.belgospatent.org.by}, 2011--\targetYear~гг.;
    \item информационно-поисковая система Федерального института промышленной собственности РФ (ФИПС) --- \url{http://www1.fips.ru}, 2011--\targetYear~гг.;
    \item всемирная база данных патентной документации Espacenet --- \url{http://ru.espacenet.com}, 2011--\targetYear~гг.;
    \item международная база данных патентной документации Google Patents --- \url{https://patents.google.com}, 2011--\targetYear~гг.
\end{itemize}

Результаты патентного поиска представлены в таблице~А.1.

\vspace{1em}
\noindent\begin{tabularx}{\textwidth}{@{}l X@{}}
Таблица А.1~--- & Результаты патентных исследований
\end{tabularx}
\vspace{0.5em}

\noindent
\begin{tabularx}{\textwidth}{|p{0.38\textwidth}|X|}
\hline
\centering Наименование аналога & \centering\arraybackslash Отличительный признак, сущность аналога \\ \hline
Пат. RU № 2420021 C2\par Способ сжатия изображений и видеопоследовательностей\par МПК G06T 9/00\par Источник: patents.google.com &
Предложен способ эффективного сжатия цветного высококачественного изображения путем пространственного разбиения и квантования компонент. Подтверждает актуальность пространственного сжатия, однако подход страдает от артефактов при высоких степенях компрессии, что в разрабатываемой системе решается применением вейвлет-декомпозиции.
\\ \hline
\end{tabularx}

\newpage

\noindent\begin{tabularx}{\textwidth}{@{}l X@{}}
Продолжение таблицы А.1
\end{tabularx}
\vspace{0.5em}

\noindent
\begin{tabularx}{\textwidth}{|p{0.38\textwidth}|X|}
\hline
\centering Наименование аналога & \centering\arraybackslash Отличительный признак, сущность аналога \\ \hline
Пат. US № 10930020 B2\par Texture compression using a neural network / Сжатие текстур с использованием нейронной сети\par МПК G06T 9/00\par Источник: patents.google.com &
Описывает метод сжатия пиксельных регионов изображений с использованием нейронной сети для генерации параметров сжатия. Решение требует значительных вычислительных мощностей на этапе кодирования, в отличие от предложенного в проекте асимметричного подхода с легковесным вейвлет-кодером.
\\ \hline
Пат. EP № 3687158 A4\par Image processing method and device / Способ и устройство обработки изображений\par МПК G06T 3/40\par Источник: patents.google.com &
Описывает метод реконструкции изображений со сверхразрешением на основе глубокой сверточной нейронной сети. Подтверждает высокую эффективность применения сверточных сетей для восстановления высокочастотных деталей изображения из сжатого представления.
\\ \hline
Пат. EP № 4258204 B1\par Image processing method, apparatus, computer program and computer-readable data carrier / Способ обработки изображений, устройство, компьютерная программа и машиночитаемый носитель данных\par МПК G06T 3/40\par Источник: patents.google.com &
Патент направлен на увеличение пространственного разрешения изображения (SISR) на основе сверточных нейросетей (CNN). Описывает архитектуру для восстановления HR-деталей из LR-изображения, что концептуально и технологически соответствует алгоритму работы декодера в разрабатываемом проекте.
\\ \hline
Пат. RU № 2741234 C1\par Способ и система для обработки изображений с использованием нейронной сети\par МПК G06T 3/40\par Источник: www1.fips.ru &
Предлагается комплексный метод нейросетевой фильтрации и масштабирования цифровых изображений.  Аналог подтверждает тенденцию выноса ресурсоемких тензорных вычислений на серверную сторону, что коррелирует с предложенной архитектурой Edge-to-Cloud.
\\ \hline
\end{tabularx}

\newpage

\noindent\begin{tabularx}{\textwidth}{@{}l X@{}}
Продолжение таблицы А.1
\end{tabularx}
\vspace{0.5em}

\noindent
\begin{tabularx}{\textwidth}{|p{0.38\textwidth}|X|}
\hline
\centering Наименование аналога & \centering\arraybackslash Отличительный признак, сущность аналога \\ \hline
Пат. US № 11006132 B2\par Image compression and decompression method, apparatus, and system using neural network / Метод, устройство и система сжатия и восстановления изображений с использованием нейросети\par МПК G06T 9/00\par Источник: patents.google.com &
Запатентована сквозная (end-to-end) архитектура нейросетевого кодирования и декодирования. Подтверждает актуальность тему, однако описанный метод требует графических ускорителей как на стороне получателя, так и на стороне отправителя, что решено в проекте путем применения асимметричного математического кодера DWT.
\\ \hline
\end{tabularx}

\vspace{1em}
Вывод: анализ патентной литературы подтверждает, что индустрия активно движется в сторону использования глубокого обучения (нейронных сетей) для задач компрессии и восстановления изображений. Выявленные аналоги (EP 3687158, EP 4258204, RU 2741234) доказывают эффективность CNN для реконструкции текстур, что обосновывает выбор моделей ESPCN и FSRCNN. В то же время патенты, описывающие исключительно нейросетевое сжатие (US 10930020, US 11006132), указывают на высокую вычислительную сложность сквозного кодирования. Предлагаемая в дипломном проекте гибридная модель (математический вейвлет-кодер Хаара + нейросетевой декодер) обладает новизной, так как решает проблему ресурсоемкости на стороне периферийного передатчика, сохраняя при этом высокое качество визуального восстановления, и не нарушает действующих патентов.

\footnotesize
\begin{lstlisting}[style=CodeListing]

    #include <iostream>
    #include <chrono>
    #include <cmath>
    #include <vector>
    #include <Windows.h>
    #include <thread>
    
    #include "include/omp.h"
    #define MOD 1000000007 
    
    long long sum(long long x, long long y)
    {
        return (x + y) % MOD;
    }
    long long mult(long long a, long long b)
    {
        return ((a % MOD) * (b % MOD)) % MOD;
    }
    long long binpow(long long x, long long n)
    {
        long long res = 1;
        while (n)
        {
            if (n & 1)
                res = (res * x) % MOD;
            x = (x * x) % MOD;
            n >>= 1;
        }
        return res;
    }
    long long divide(long long a, long long b)
    {
        //return mult(a, binpow(b, phi(MOD) - 1)) % MOD;
        return mult(a, binpow(b, MOD - 2)) % MOD;
    }
    
    void calculate_sum(long long tmax)
    {
        long long taylor = 0;
        long long x = 17, nx = 1, fa = 1;
        #pragma omp parallel for
        for(int i = 0; i < tmax; i++)
        {
            if(i > 0) fa = mult(fa, i);
            nx = mult(nx, x);
            taylor = sum(taylor, divide(nx, fa));
        }
    }
    
    std::vector<long long> v;
    std::vector<long long> table;
    
    const unsigned long long MAX = 1000000;
    const long double stepSize = 0.02d;
    const long long tableStepSize = 10;
    
    int main()
    {
        // freopen("input.txt", "r", stdin);
        // freopen("output_NO_OMP.txt", "w", stdout);
        // freopen("outputOMP.txt", "w", stdout);
    
        MEMORYSTATUSEX status;
        status.dwLength = sizeof(status);
        GlobalMemoryStatusEx(&status);
        unsigned long long totalMemory = status.ullTotalPhys;
        double totalMemoryGB = static_cast<double>(totalMemory) / (1024 * 1024 * 1024);
        std::cout << "Total physical memory: " << totalMemoryGB << " GB" << std::endl;
    
        int numThreads = omp_get_max_threads();
        std::cout << "All threads : " << numThreads << std::endl;
        
        omp_set_num_threads(1);
        std::cout << "1 thread : " << std::endl;
        auto begin = std::chrono::high_resolution_clock::now();
        for(long long tmax = MAX * stepSize; tmax <= MAX; tmax += MAX * stepSize)
        {
            calculate_sum(tmax);
            auto end = std::chrono::high_resolution_clock::now();
            auto duration = std::chrono::duration_cast<std::chrono::microseconds>(end - begin);
            v.push_back(duration.count());
        }
        #pragma omp master
        {
            for(int i = 0; i < v.size(); i += 1)
                std::cout << v[i] << ", ";
        }
        std::cout << std::endl;
        for(size_t i = 0; i < v.size(); i += tableStepSize)
            table.push_back(v[i + 9]);
        v.clear();
        
        // Измерение времени выполнения на всех физических потоках
        omp_set_num_threads(numThreads / 2);
        std::cout << numThreads / 2 << " threads (all physical threads): " << std::endl;
        begin = std::chrono::high_resolution_clock::now();
        for(long long tmax = MAX * stepSize; tmax <= MAX; tmax += MAX * stepSize)
        {
            calculate_sum(tmax);
            auto end = std::chrono::high_resolution_clock::now();
            auto duration = std::chrono::duration_cast<std::chrono::microseconds>(end - begin);
            v.push_back(duration.count());
        }
        #pragma omp master
        {
            for(int i = 0; i < v.size(); i += 1)
                std::cout << v[i] << ", ";
        }
        std::cout << std::endl;
        for(size_t i = 0; i < v.size(); i += tableStepSize)
            table.push_back(v[i + 9]);
        v.clear();
    
        // Измерение времени выполнения на всех логических потоках
        omp_set_num_threads(numThreads);
        std::cout << numThreads << " threads (all logical threads): " << std::endl;
        begin = std::chrono::high_resolution_clock::now();
        for(long long tmax = MAX * stepSize; tmax <= MAX; tmax += MAX * stepSize)
        {
            calculate_sum(tmax);
            auto end = std::chrono::high_resolution_clock::now();
            auto duration = std::chrono::duration_cast<std::chrono::microseconds>(end - begin);
            v.push_back(duration.count());
        }
        #pragma omp master
        {
            for(int i = 0; i < v.size(); i += 1)
                std::cout << v[i] << ", ";
        }
        std::cout << std::endl;
        for(size_t i = 0; i < v.size(); i += tableStepSize)
            table.push_back(v[i + 9]);
        v.clear();
    
        std::cout << "Execution time comparison:" << std::endl;
        std::cout << "-------------------------------------------------------------------------" << std::endl;
        std::cout << "| Threads\t| Time (in milliseconds)\t\t\t\t|" << std::endl;
        std::cout << "-------------------------------------------------------------------------" << std::endl;
        std::cout << "| 1\t\t| ";
        for(size_t i = 0; i < 5; i++)
            std::cout << table[i] << " ";
        std::cout << "\t\t|" << std::endl;
        std::cout << "| Half\t\t| ";
        for(size_t i = 5; i < 10; i++)
            std::cout << table[i] << " ";
        std::cout << "\t\t\t|" << std::endl;
        std::cout << "| All \t\t| ";
        for(size_t i = 10; i < 15; i++)
            std::cout << table[i] << " ";
        std::cout << "\t\t\t|" << std::endl;
        std::cout << "-------------------------------------------------------------------------" << std::endl;
        return 0;
    }

\end{lstlisting}
\normalsize